\chapter{Ablation Studies and Scientific Evaluation}
\label{chap:ch21}

\section{The Scientific Method in Deep Learning Research}
In academic machine learning, asserting that an architecture is "better" is meaningless without rigorous, controlled \textbf{ablation studies}. An ablation study systematically isolates each proposed component to prove whether performance gains stem from genuine architectural innovations or random statistical noise.

\section{Systematic Experimental Evolution Across Four Stages}
Table~\ref{tab:ablation_results} summarizes the quantitative progression across the four major experimental iterations conducted in our repository, evaluated across held-out heavy-cloud validation scenes ($>70\%$ cloud cover).

\begin{table}[htbp]
\centering
\small
\begin{tabularx}{\textwidth}{lccccX}
\toprule
\textbf{Experimental Stage} & \textbf{Heavy PSNR (dB)} & \textbf{Heavy SSIM} & \textbf{Detail Ratio} & \textbf{Edge Corr} & \textbf{Reconstruction Character} \\
\midrule
\textbf{Stage 1: Baseline ECRformer} & \textbf{30.20} & 0.8747 & 0.3540 & 0.2812 & Severe over-smoothing; flat, painted appearance; $65\%$ texture lost. \\
\textbf{Stage 2: Spectral FFT Loss} & 29.85 & 0.8690 & 0.5120 & 0.3140 & Moderate texture recovery; periodic grid artifacts from Fourier ringing. \\
\textbf{Stage 3: Structural/Perceptual} & 30.08 & 0.8792 & 0.6845 & 0.3725 & Sharp boundary contours; strong vegetation contrast; realistic macroscopic layout. \\
\textbf{Stage 4: Sweetspot PatchGAN} & 29.97 & \textbf{0.8815} & \textbf{0.8576} & \textbf{0.4153} & Near-total texture preservation ($>85\%$); authentic surface grain; sharp roads. \\
\bottomrule
\end{tabularx}
\caption{Systematic ablation and benchmark results evaluated across held-out validation patches with heavy cloud coverage ($>70\%$).}
\label{tab:ablation_results}
\end{table}

\section{In-Depth Analysis of Ablation Results}

\subsection{Stage 1 vs. Stage 4: Deconstructing the "PSNR Trap"}
Notice that Stage 1 achieved the highest raw PSNR ($30.20\text{ dB}$), yet suffered from catastrophic texture loss (Detail Ratio $0.3540$). 
Because Stage 1 minimized pure $L_1$ error, the model learned to output smooth, safe conditional averages. 

In Stage 4, adding the 11-channel high-pass PatchGAN and local standard deviation loss caused the Detail Ratio to surge from $\mathbf{0.3540 \to 0.8576}$---a $\mathbf{+142\%}$ improvement in texture retention---with SSIM climbing from $0.8747 \to 0.8815$. The minor $-0.23\text{ dB}$ shift in PSNR is a classic manifestation of the perception-distortion trade-off: slight sub-pixel phase adjustments yield vastly superior visual realism and physical fidelity!

\subsection{The Impact of Intermediate SDFL Auxiliary Supervision}
Ablating the auxiliary projection heads (\texttt{down\_proj[0..2]} and \texttt{up\_proj[0..2]}) showed that training took $2.4\times$ longer to reach equivalent validation loss. Without multi-scale supervision, early encoder layers experienced sluggish gradient updates during the first 10 epochs.

\subsection{The Impact of Soft-Mask Blending}
Evaluating predictions without the soft-mask blend gate---using the raw neural output $\mathbf{P}_{\text{dir}}$ globally---caused PSNR on clear-sky test patches to drop from $38.4\text{ dB} \to 28.1\text{ dB}$. Even subtle neural imperfections degrade pristine optical observations. Soft-mask blending is mandatory to ensure clear pixels remain $100\%$ untouched.

\section{Computational Efficiency and Parameter Footprint}
\begin{itemize}
    \item \textbf{Active Generator Parameter Count:} 3,619,070 parameters ($3.62\text{ M}$).
    \item \textbf{Disk Checkpoint Size:} $\approx 34.3\text{ MB}$ (containing generator, EMA shadow weights, and PatchGAN discriminator).
    \item \textbf{Inference Speed on RTX 5070:} $<12\text{ ms}$ per $128 \times 128$ patch ($>80\text{ patches/second}$).
    \item \textbf{Mobile / Edge Feasibility:} The standalone generator can be converted to LiteRT / ONNX ($\approx 14.5\text{ MB}$ in float32, or $<4.0\text{ MB}$ in INT8 quantization), enabling real-time deployment on edge devices or Android ground stations!
\end{itemize}

\section{Summary of Chapter 21}
\begin{enumerate}
    \item Controlled ablation benchmarks demonstrate how each architectural modification directly targets specific physical deficiencies.
    \item Moving from Stage 1 to Stage 4 improved high-frequency detail retention from $35.4\%$ to $85.8\%$.
    \item The perception-distortion trade-off explains why human-preferred, photorealistic reconstructions can have marginally lower PSNR than artificially blurred predictions.
    \item Auxiliary SDFL heads accelerate convergence by $2.4\times$.
    \item The active 3.62M generator is computationally lightweight, fast, and primed for edge deployment.
\end{enumerate}

\section{Practice Exercises}
\begin{enumerate}
    \item Why is a higher PSNR not always proof that a cloud removal model is superior?
    \item What quantitative metric in Table~\ref{tab:ablation_results} best captures the transition from a blurry reconstruction to a sharp reconstruction?
    \item Why does omitting the soft-mask blend gate significantly degrade clear-sky test scores?
\end{enumerate}

\subsection*{Solutions}
\begin{enumerate}
    \item PSNR measures logarithmic pixel-wise squared error. An over-smoothed model that predicts the uniform conditional mean avoids extreme pixel penalties and can score a higher PSNR than a model that synthesizes sharp, authentic textures with slight sub-pixel phase offsets.
    \item The \textbf{Detail Ratio} (Laplacian high-frequency energy retention) best captures texture sharpness, rising from $0.3540$ in the blurry baseline to $0.8576$ in the sharp GAN model.
    \item The original clear optical pixels have zero error relative to the clear ground truth. Passing clear terrain through a generative neural network invariably introduces minor neural noise and color shifts, whereas soft-mask blending copies the true optical pixels with zero loss of fidelity.
\end{enumerate}
